\section{Gap Taxonomy: Four Types, Four Zones}
\label{sec:taxonomy}

We identify four geographic gap zones and classify gap counties within each zone by type.

\subsection{Zone 1 — The Northern Tier (Type 1: T1 Geographic Gap)}

\textbf{Extent}: The strip of counties along the US–Canada border from Minnesota through Montana, corresponding to the northern halves of Minnesota, North Dakota, and Montana, plus the Northern Michigan Upper Peninsula.

\textbf{Gap counties}: 44 counties in Montana (502,001 people), 44 in North Dakota (324,862 people), 50 in Minnesota (1,283,216 people, including parts of the Iron Range), and 43 in Michigan (1,388,641 people, Upper Peninsula).

\textbf{Nearest existing interstate}: I-90 (southern Minnesota/Wyoming/Montana) is the closest T1 corridor. The I-90/Canadian border gap is 300--500 miles across most of Montana and North Dakota.

\textbf{Cause}: The original interstate system connected major population centers of the 1950s. The northern tier --- predominantly agricultural, with sparse population --- was not prioritized. US-2, the primary east-west route through this region, runs 2,571 miles from Houlton, ME to Everett, WA at two-lane highway standard, with no interstate segments.

\textbf{Proposed resolution}: A Northern Tier interstate along the US-2 alignment would serve the contiguous northern Montana/North Dakota/Minnesota gap zone directly. It would be the longest proposed interstate in the US at $\sim$2,500 miles, qualifying as a T1 corridor under the centrality-adjusted classification of \citet{ROUTE_A1} given its irreplaceability (B1 score $\geq 9.5$; no parallel route within 300 miles).

\textbf{Type classification}: Type 1 (T1 Geographic Gap). The entire zone lacks a T1 corridor. The Upper Michigan gap (Keweenaw, Luce, Houghton, Ontonagon counties) is Type 3 (Rural Isolation) due to the Upper Peninsula's terrain and low population density; a T2 connector (I-41 extension or US-2 upgrade through the UP) would be more appropriate than a full T1 designation.

\subsection{Zone 2 — The Appalachians (Type 2: T2 Missing Links)}

\textbf{Extent}: West Virginia (26 counties, 517,756 people), eastern Kentucky (66 counties, 1,698,632 people), western Virginia (30 counties, 657,852 people), and eastern Tennessee (42 counties, 1,660,090 people).

\textbf{Nearest existing interstates}: I-64/I-77/I-79 serve the western fringe of West Virginia; I-81 serves the western Virginia valley; I-40 passes through Knoxville. But no interstate crosses the central Appalachian ridge-and-valley system from east to west between I-40 (Knoxville) and I-70 (Wheeling, WV) --- a 400-mile gap.

\textbf{Cause}: The Appalachian Mountains present severe engineering challenges for interstate construction. All interstate crossings of the Appalachians occur at lower elevations (I-40 through the Tennessee gaps, I-70 through the National Road corridor, I-68 through a Maryland notch). The central Appalachian spine in WV/VA/KY has no comparable passage.

\textbf{Economic dimension}: Appalachian gap counties have some of the highest C3 scores in the dataset --- mean 7.2, corresponding to buffer GDP per capita at approximately 65\% of the national average. This is the region identified in USDA ``deep rural'' analyses as having the worst outcomes across economic, health, and educational dimensions \citep{USDA_ERS2020}. The Appalachian Regional Commission has documented this deficit for 60 years \citep{ARC2023}; the missing interstate link is one structural contributor.

\textbf{Proposed resolution}: The I-3 corridor (proposed Savannah, GA to Detroit, MI via Knoxville, Lexington, and Toledo) would cross the Appalachians through existing river valleys and provide the first E-W interstate through the central Appalachian region. An Appalachian connector (West Virginia Turnpike extension upgraded to full interstate standard) would provide a shorter alternative.

\textbf{Type classification}: Primarily Type 2 (T2 Missing Links) for the WV/KY/VA interior counties. A full T1 designation is not warranted for most of this zone given the parallel I-64/I-81 options; what's needed is a T2 E-W connector that reduces detour distances. The eastern Kentucky coalfield counties are Type 4 (Economic Opportunity Gaps) --- both geographically isolated and at the bottom of national economic rankings.

\subsection{Zone 3 — Rural Gulf South (Type 2 and Type 4)}

\textbf{Extent}: Louisiana (45 parishes, 2,631,301 people), Mississippi (40 counties, 752,314 people), Arkansas (50 counties, 1,116,108 people), and rural Texas (145 counties, 3,459,424 people --- the largest gap population in any state).

\textbf{Nearest existing interstates}: I-10 runs the Gulf Coast; I-20 runs 150--200 miles north. The zone between them across Louisiana, Mississippi, and Arkansas has limited interstate coverage. Rural Texas lacks interstate access across a vast triangle formed by I-10, I-20, and I-35.

\textbf{Cause}: The Gulf South interior was served by US highways (US-90, US-82, US-190) that were never upgraded to interstate standard. The delta regions (Mississippi Delta, Atchafalaya Basin) have challenging terrain for interstate construction. The Texas triangle interior (between San Antonio, Fort Worth, and Houston) is largely agricultural and was not prioritized in the original system.

\textbf{Proposed resolution}: I-14 (partially designated in Texas as a new east-west corridor) would serve the Gulf Coast interior. I-69 completion through East Texas would bring rural Texas counties within range of the HOU-CHI corridor. The Meridian corridor (I-20 to Gulf connecting Mississippi) has been proposed in state LRTPs.

\textbf{Type classification}: Mixed Type 2 (T2 Missing Links for higher-population parishes/counties that would be served by I-14 or I-69) and Type 4 (Economic Opportunity Gaps for the Mississippi Delta and rural Appalachian portions of Arkansas, where C3 scores average 7.8).

\subsection{Zone 4 — Rural West (Type 3: Rural Isolation)}

\textbf{Extent}: New Mexico (21 counties, 874,061 people), Nevada (13 counties, 727,071 people), Utah (20 counties, 1,353,166 people), Colorado (44 counties, 1,112,493 people), Idaho (28 counties, 598,980 people).

\textbf{Nearest existing interstates}: I-40 and I-10 in New Mexico; I-15 and I-80 in Nevada/Utah; I-70, I-25, I-76 in Colorado. But the vast geographic scale of western counties means county centroids can be hundreds of miles from any on-ramp even when the county itself contains an interstate segment.

\textbf{Cause}: Western states have counties of enormous geographic scale (San Bernardino County, CA: 20,000 sq mi; Nye County, NV: 18,182 sq mi) where population is concentrated in a county seat far from the county centroid. This is partly a county-centroid artifact. However, it also reflects genuine access deficits: communities like Winnemucca, NV; Ely, NV; and Alamogordo, NM are genuinely isolated from the interstate system.

\textbf{Note on California}: California's 34 gap counties and 12.2 million gap population is dominated by the county-centroid effect for large inland counties. Los Angeles County itself has its centroid in the Santa Monica Mountains, appearing as a gap county despite containing over 100 miles of interstate. San Bernardino County's centroid falls in the Mojave Desert 80 miles from the nearest I-10 on-ramp. We flag California gap counties for separate analysis; the county-centroid measure overstates the California coverage deficit significantly.

\textbf{Proposed resolution}: Rural access spurs (Section~\ref{sec:corridors}) rather than new full interstates for most Type 3 counties. Isolated communities with genuine distance from the interstate (not just centroid artifact) are candidates for T3 rural connector spurs of $\leq 10$ miles.

\textbf{Type classification}: Primarily Type 3 (Rural Isolation) --- low population, large land area, difficult terrain. Few proposed corridor solutions exist and the coverage efficiency metric is unfavorable (few people served per dollar of new construction in these counties). Rural access spurs at strategic points are the appropriate investment.

\subsection{Nationwide Summary}

\begin{table}[h!]
\centering
\caption{Gap Zone Summary by Type and Population}
\label{tab:zones}
\begin{tabular}{llrr}
\toprule
\textbf{Zone} & \textbf{Primary Type} & \textbf{Gap Counties} & \textbf{Gap Population} \\
\midrule
Northern Tier (ND/MT/MN/MI UP) & T1 Geographic Gap & $\sim$181 & $\sim$3,499,000 \\
Appalachians (WV/KY/VA/TN) & T2 Missing Links + T4 Eco & $\sim$164 & $\sim$4,534,000 \\
Rural Gulf South (LA/MS/AR/TX) & T2 Missing Links + T4 Eco & $\sim$280 & $\sim$7,959,000 \\
Rural West (NM/NV/UT/CO/ID/CA) & T3 Rural Isolation & $\sim$160 & $\sim$16,840,000 \\
All other gap counties & Various & $\sim$725 & $\sim$33,681,000 \\
\midrule
\textbf{Total continental US} & & \textbf{1,510} & \textbf{66,513,310} \\
\bottomrule
\end{tabular}
\end{table}
